\documentclass[11pt]{article}

% ---------- Packages ----------
\usepackage[utf8]{inputenc}
\usepackage[T1]{fontenc}
\usepackage{amsmath}
\usepackage{amssymb}
\usepackage[margin=1in]{geometry}
\usepackage{titlesec}
\usepackage{enumitem}
\usepackage{xcolor}
\usepackage{longtable}
\usepackage{booktabs}
\usepackage{array}
\usepackage{parskip}
\usepackage{hyperref}

% ---------- Colors ----------
\definecolor{primary}{HTML}{1B3A5C}
\definecolor{accent}{HTML}{9C3D1B}
\definecolor{lightgray}{HTML}{5A5A5A}

\hypersetup{
  colorlinks=true,
  linkcolor=primary,
  urlcolor=accent,
  citecolor=primary,
  pdftitle={Course Breakdown: Calculus for the Biological and Social Sciences},
  pdfauthor={Course Design Summary}
}

% ---------- Heading styles ----------
\titleformat{\section}
  {\normalfont\Large\bfseries\color{primary}}{\thesection}{0.8em}{}[{\color{primary!40}\titlerule}]
\titleformat{\subsection}
  {\normalfont\large\bfseries\color{primary}}{\thesubsection}{0.8em}{}
\titleformat{\subsubsection}
  {\normalfont\normalsize\bfseries\color{accent}}{}{0em}{}
\titlespacing*{\subsubsection}{0pt}{1.4em}{0.6em}

\titleformat{\paragraph}
  {\normalfont\bfseries\itshape\color{lightgray}}{}{0em}{}
\titlespacing*{\paragraph}{0pt}{0.9em}{0.3em}

\setcounter{secnumdepth}{2}
\setcounter{tocdepth}{3}

\setlist[itemize]{itemsep=1pt, topsep=2pt, parsep=0pt, leftmargin=1.4em}
\setlist[enumerate]{itemsep=2pt, topsep=2pt, parsep=0pt, leftmargin=1.6em}

% Convenience commands for the recurring five-part chapter template
\newcommand{\ChapHead}[2]{\subsubsection{#1 \textnormal{\normalfont\normalsize\color{lightgray}(pp.\ #2)}}}
\newcommand{\MajorTopics}{\paragraph{Major Topics}}
\newcommand{\KeyConcepts}{\paragraph{Key Concepts}}
\newcommand{\LearnObj}{\paragraph{Learning Objectives --- students will be able to:}}
\newcommand{\PreReq}{\paragraph{Prerequisite Knowledge}}
\newcommand{\Hard}{\paragraph{Challenging Topics \& Teaching Notes}}

\title{\color{primary}\textbf{Course Breakdown}\\[0.3em]
\Large Calculus for the Biological and Social Sciences}
\author{Derived from the course textbook's Table of Contents}
\date{\today}

\begin{document}

\maketitle

\begin{center}
\begin{minipage}{0.85\textwidth}
\small\itshape\centering
A structured curriculum map --- course structure, modules, chapters, key concepts,
learning objectives, prerequisites, and challenging topics --- built from the
textbook's table of contents for a life-science- and social-science-oriented
Calculus sequence.
\end{minipage}
\end{center}

\tableofcontents
\newpage

% =====================================================================
\section*{0. Course Overview}
\addcontentsline{toc}{section}{0. Course Overview}

The source textbook covers twelve chapters that blend standard single- and
multivariable calculus with discrete-time modeling, differential equations,
linear algebra, and probability/statistics --- the toolkit most commonly used
in biology and social-science quantitative courses. Twelve chapters is more
material than a single semester of ``Calculus 1'' typically covers, so this
breakdown makes the following working assumption, stated here explicitly:

\begin{quote}\small
\textbf{Assumption.} The textbook is delivered across a two-semester sequence.
\textbf{Calculus I} (this course) corresponds to \textbf{Chapters 1--6}
(precalculus review through integration). \textbf{Calculus II} corresponds to
\textbf{Chapters 7--12} (integration techniques through probability and
statistics) and is included here so the breakdown reflects the complete table
of contents and shows how Calculus I feeds into it. If your program instead
runs this material as a single accelerated semester, a three-quarter sequence,
or omits later chapters, the six modules below can be regrouped or trimmed
without changing their internal structure.
\end{quote}

\paragraph{Audience and emphasis} Students are primarily biology, ecology,
pre-health, and social-science majors. Motivating examples throughout the
course are drawn from population growth, drug kinetics, epidemiology,
ecological interactions, demography, and data analysis rather than physics.

\paragraph{Asterisk convention} Following the textbook's own convention,
topics marked with an asterisk (*) are optional or enrichment material --
deeper rigor (e.g.\ formal $\varepsilon$--$\delta$ limits), extensions of core
ideas (e.g.\ stability of recurrence equations), or additional modeling
examples. A course can be taught in full without covering starred topics;
instructors choose which to include based on time and student background.

\paragraph{Computational component} Several sections (numerical recursion,
the bisection method, numerical integration) are built around spreadsheet
use. A weekly computer lab or homework component using a spreadsheet
application is recommended alongside lecture.

% =====================================================================
\section*{1. Overall Prerequisites (Before Calculus I)}
\addcontentsline{toc}{section}{1. Overall Prerequisites (Before Calculus I)}

\begin{itemize}
\item \textbf{Algebra:} manipulating and solving linear, polynomial, and
rational equations and inequalities; working with exponents and radicals.
\item \textbf{Functions and graphs:} function notation, domain and range,
composition, and basic graph reading.
\item \textbf{Geometry:} equations of lines and circles; basic coordinate
geometry.
\item \textbf{Trigonometry:} the unit circle, radian measure, and core
identities.
\item \textbf{Exponentials and logarithms:} laws of exponents and logarithms,
and the relationship between them.
\item \textbf{Complex numbers:} arithmetic of complex numbers and the
quadratic formula for non-real roots.
\end{itemize}
Chapter 1 of the course is itself a diagnostic-and-review chapter, so gaps in
the list above are expected to be identified and remediated in the first two
weeks rather than assumed away.

% =====================================================================
\section*{2. Course Structure at a Glance}
\addcontentsline{toc}{section}{2. Course Structure at a Glance}

The twelve chapters are grouped into six two-chapter modules, three per
semester.

\begin{longtable}{@{}p{0.09\textwidth}p{0.30\textwidth}p{0.10\textwidth}p{0.10\textwidth}p{0.32\textwidth}@{}}
\toprule
\textbf{Module} & \textbf{Chapters} & \textbf{Semester} & \textbf{Approx.\ Weeks} & \textbf{Theme} \\
\midrule
\endhead
A & 1. Preview and Review \newline 2. Discrete-Time Models & Calc.\ I & 4 & Foundations: precalculus review and discrete-time (recurrence-equation) modeling \\
B & 3. Limits and Continuity \newline 4. Differentiation & Calc.\ I & 5 & Limits as the gateway to the derivative; full toolkit of differentiation rules \\
C & 5. Applications of Differentiation \newline 6. Integration & Calc.\ I & 6 & Using derivatives to analyze functions; the definite integral and the Fundamental Theorem \\
D & 7. Integration Techniques \newline 8. Differential Equations & Calc.\ II & 5 & Symbolic/numerical integration methods; first-order differential equation models \\
E & 9. Linear Algebra \& Analytic Geometry \newline 10. Multivariable Calculus & Calc.\ II & 5 & Matrices, eigenvalues, and calculus of several variables \\
F & 11. Systems of Differential Equations \newline 12. Probability and Statistics & Calc.\ II & 5 & Coupled dynamical systems and the language of probability/statistics \\
\bottomrule
\end{longtable}

% =====================================================================
\section{Semester 1: Calculus I}

\subsection{Module A --- Foundations: Functions and Discrete-Time Models}

\ChapHead{Chapter 1: Preview and Review}{1--61}

\MajorTopics
\begin{itemize}
\item Precalculus skills diagnostic test
\item Preliminaries: real numbers, lines, circles, trigonometry, exponentials/logarithms, complex numbers and quadratics
\item Elementary functions: polynomial, rational, power, exponential, inverse, logarithmic, trigonometric
\item Graphing: transformations, logarithmic scales, transforming data into linear form, (optional) verbal description to graph
\end{itemize}

\KeyConcepts
Function families and their defining properties; domain/range; graph
transformations (shifts, stretches, reflections); inverse functions;
logarithmic and semi-log axes as a way to linearize exponential or
power-law data.

\LearnObj
\begin{enumerate}
\item Use the diagnostic test to identify and remediate personal precalculus gaps.
\item Solve equations involving lines, circles, exponentials, logarithms, and quadratics (including non-real roots).
\item Classify a function by family and describe its domain, range, and key graphical features.
\item Apply shifts, stretches, and reflections to graph a function and its inverse.
\item Re-express exponential or power-law data on logarithmic or semi-log axes to reveal a linear trend.
\item * Translate a verbal description of a process into a qualitative graph.
\end{enumerate}

\PreReq High-school algebra, geometry, and trigonometry (see Section 1 above).

\Hard Complex numbers and the quadratic formula tend to be rusty for
returning students; log and semi-log scales are conceptually new even to
students comfortable with logarithms as a computation. Building a
qualitative graph directly from a verbal story (Sec.\ 1.4.4*) is a skill,
not a formula, and benefits from in-class modeling.

\bigskip
\ChapHead{Chapter 2: Discrete-Time Models, Sequences, and Difference Equations}{62--100}

\MajorTopics
\begin{itemize}
\item Exponential growth and decay: discrete-time population models, recurrence equations, visualization
\item Sequences: definitions, (optional) spreadsheet calculation, limits, recurrence equations, sigma notation
\item Modeling with recurrence equations: density-dependent growth, the Beverton--Holt model, the discrete logistic equation, (optional) drug absorption
\end{itemize}

\KeyConcepts
Recurrence (difference) equations; discrete exponential growth
$N_{t+1}=RN_t$; cobweb/stair-step diagrams; convergence of a sequence;
sigma notation; density dependence and carrying capacity; the
Beverton--Holt and discrete logistic maps.

\LearnObj
\begin{enumerate}
\item Formulate a recurrence equation that models discrete-time population growth or decay.
\item Iterate and visualize solutions of a recurrence equation using cobweb/stair-step diagrams.
\item Determine whether a sequence converges and, if so, find its limit.
\item Evaluate sums of sequences using sigma notation.
\item Compare density-independent and density-dependent growth models and interpret carrying capacity.
\item * Model drug absorption using a discrete recurrence equation.
\end{enumerate}

\PreReq Exponential and logarithmic functions (Sec.\ 1.3.5, 1.3.7); comfort with algebraic manipulation of equations.

\Hard Predicting the long-term (equilibrium or oscillatory) behavior of a
\emph{nonlinear} recurrence equation is the first genuinely modeling-heavy
task in the course and is often harder than the algebra itself; students
also frequently struggle to set up a recurrence equation directly from a
word problem rather than recognizing a pre-built formula.

\subsection{Module B --- Limits and Differential Calculus}

\ChapHead{Chapter 3: Limits and Continuity}{101--141}

\MajorTopics
\begin{itemize}
\item Limits: an informal treatment, common pitfalls, and limit laws
\item Continuity: definition and combinations of continuous functions
\item Limits at infinity
\item Trigonometric limits and the Sandwich (Squeeze) Theorem
\item Properties of continuous functions: the Intermediate-Value Theorem and the bisection method
\item * A formal ($\varepsilon$--$\delta$) definition of limits
\end{itemize}

\KeyConcepts
One-sided and two-sided limits; limit laws; continuity at a point and on an
interval; removable, jump, and infinite discontinuities; horizontal and
vertical asymptotes; the Sandwich Theorem; the Intermediate-Value Theorem
(IVT); the bisection method for root-finding.

\LearnObj
\begin{enumerate}
\item Evaluate limits using direct substitution, algebraic manipulation, and limit laws.
\item Identify discontinuities of a function and classify their type.
\item Evaluate limits at infinity and identify horizontal and vertical asymptotes.
\item Apply the Sandwich Theorem to evaluate trigonometric limits.
\item Use the IVT to justify the existence of a root and approximate it with the bisection method.
\item * State and apply the formal $\varepsilon$--$\delta$ definition of a limit.
\end{enumerate}

\PreReq Elementary functions and their graphs (Chapter 1); the intuitive notion of a sequence limit (Sec.\ 2.2.3).

\Hard The formal $\varepsilon$--$\delta$ definition (Sec.\ 3.6*) is the
first rigor-heavy topic in the course and is genuinely difficult for a
non-math-major audience; trigonometric limits requiring the Sandwich
Theorem require comfort manipulating inequalities, and students often
confuse ``the limit does not exist'' with ``the function is undefined.''

\bigskip
\ChapHead{Chapter 4: Differentiation}{142--212}

\MajorTopics
\begin{itemize}
\item Formal definition of the derivative
\item Properties of the derivative: interpretation; differentiability vs.\ continuity
\item The power rule and derivatives of polynomials
\item The product and quotient rules; derivatives of rational and power functions
\item The chain rule (with proof)
\item Implicit functions, implicit differentiation, and related rates
\item Higher derivatives
\item Derivatives of trigonometric, exponential, inverse, and logarithmic functions; the inverse tangent function
\item Linear approximation and error propagation
\end{itemize}

\KeyConcepts
The derivative as a limit of a difference quotient and as an instantaneous
rate of change / slope of the tangent line; differentiability implies
continuity (but not conversely); the power, product, quotient, and chain
rules; implicit differentiation; related rates; higher-order derivatives;
linear approximation and differentials.

\LearnObj
\begin{enumerate}
\item Compute a derivative directly from its limit definition.
\item Apply the power, product, quotient, and chain rules to differentiate combinations of functions.
\item Perform implicit differentiation and solve related-rates problems.
\item Differentiate trigonometric, exponential, logarithmic, and inverse functions.
\item Compute higher-order derivatives.
\item Use linear approximation and differentials to estimate function values and propagate measurement error.
\end{enumerate}

\PreReq Limits and continuity (Chapter 3); the elementary function families of Chapter 1.

\Hard Related-rates and implicit-differentiation word problems are
consistently among the hardest topics in introductory calculus because they
require translating a physical/biological scenario into an equation before
any calculus is applied; combining the chain rule with implicit
differentiation compounds this. Error propagation (Sec.\ 4.11) is a valuable
but unfamiliar application for many students.

\subsection{Module C --- Applications of Differentiation and Integral Calculus}

\ChapHead{Chapter 5: Applications of Differentiation}{213--305}

\MajorTopics
\begin{itemize}
\item Extrema and the Mean-Value Theorem
\item Monotonicity and concavity
\item Extrema and inflection points
\item Optimization
\item L'Hopital's Rule
\item Graphing and asymptotes
\item * Stability of recurrence equations
\item * The Newton--Raphson method
\item * Modeling biological systems with differential equations
\item Antiderivatives
\end{itemize}

\KeyConcepts
Critical points; local vs.\ global extrema; the Extreme-Value Theorem and
the Mean-Value Theorem; the first- and second-derivative tests;
concavity and inflection points; constrained optimization; indeterminate
forms; curve sketching; equilibrium stability; Newton--Raphson iteration;
antiderivatives.

\LearnObj
\begin{enumerate}
\item Locate and classify local and global extrema using the first- and second-derivative tests.
\item Apply and interpret the Mean-Value Theorem.
\item Solve applied optimization problems (e.g.\ resource allocation, foraging efficiency).
\item Evaluate indeterminate-form limits using L'Hopital's Rule.
\item Sketch a complete graph of a function, including asymptotes, extrema, and inflection points.
\item Compute basic antiderivatives.
\item * Determine the stability of equilibria of a discrete recurrence model; apply Newton--Raphson to approximate roots; set up an introductory differential-equation model of population growth.
\end{enumerate}

\PreReq Differentiation rules (Chapter 4); recurrence equations (Chapter 2, for the starred stability section).

\Hard Translating a word problem into an optimization objective function
and constraint(s) is the single hardest recurring skill in this chapter;
equilibrium/stability analysis of a nonlinear recurrence equation (Sec.\
5.7*) revisits Chapter 2 with new tools and is conceptually demanding;
students often misapply L'Hopital's Rule to forms that are not actually
indeterminate.

\bigskip
\ChapHead{Chapter 6: Integration}{306--354}

\MajorTopics
\begin{itemize}
\item The definite integral: the area problem, the general theory of Riemann integrals, properties (including optional order properties)
\item The Fundamental Theorem of Calculus, Parts I and II (with an optional rigorous treatment via Leibniz's Rule); antiderivatives and indefinite integrals
\item Applications: cumulative change, average values, and (optional) the Mean-Value Theorem for integrals, areas, volumes of solids, and rectification of curves
\end{itemize}

\KeyConcepts
Riemann sums; the definite integral as a limit of Riemann sums / signed
area / accumulated change; linearity and additivity of the integral; the
Fundamental Theorem of Calculus (both parts); the indefinite integral;
average value of a function; area between curves; volumes of revolution;
arc length.

\LearnObj
\begin{enumerate}
\item Construct and evaluate Riemann sums and interpret the definite integral as their limit.
\item Apply properties of the definite integral (linearity, additivity, comparison).
\item Apply the Fundamental Theorem of Calculus to evaluate definite integrals and to differentiate accumulation functions.
\item Compute indefinite integrals of standard functions.
\item Use integration to compute cumulative change and average values (e.g.\ total drug dose, average population size).
\item * Compute areas between curves, volumes of solids of revolution, and arc length.
\end{enumerate}

\PreReq Antiderivatives (Sec.\ 5.10); limits (Chapter 3, for the Riemann-sum limiting process).

\Hard The rigorous construction of the Riemann integral (Sec.\ 6.1.2) and
the proof of the Fundamental Theorem via Leibniz's Rule (Sec.\ 6.2.2*) are
the most abstract material so far; volume and arc-length integrals (Sec.\
6.3.5*--6.3.6*) require strong three-dimensional/geometric visualization
that many students have not practiced since high-school geometry.

% =====================================================================
\newpage
\section{Semester 2: Calculus II}

\subsection{Module D --- Integration Techniques and Differential Equations}

\ChapHead{Chapter 7: Integration Techniques and Computational Methods}{355--426}

\MajorTopics
\begin{itemize}
\item The substitution rule (indefinite and definite integrals)
\item Integration by parts and practicing integration
\item Rational functions and partial fractions (proper rational functions, decomposition, repeated linear factors, (optional) irreducible quadratic factors, summary)
\item * Improper integrals (unbounded intervals, unbounded integrand, a comparison result)
\item Numerical integration: the midpoint rule, the trapezoidal rule, spreadsheet implementation, (optional) error estimation
\item * The Taylor approximation
\item * Tables of integrals
\end{itemize}

\KeyConcepts
$u$-substitution; integration by parts; partial-fraction decomposition;
improper integrals and convergence/divergence; the comparison test;
midpoint and trapezoidal numerical rules; Taylor polynomials and their
remainder/error.

\LearnObj
\begin{enumerate}
\item Evaluate integrals using $u$-substitution and integration by parts.
\item Decompose a rational function into partial fractions (including repeated and, optionally, irreducible quadratic factors) and integrate it.
\item * Determine convergence or divergence of an improper integral and evaluate convergent ones.
\item Approximate a definite integral numerically using the midpoint and trapezoidal rules, and estimate the error.
\item * Construct a Taylor polynomial to approximate a function and bound the approximation error.
\end{enumerate}

\PreReq The definite/indefinite integral and the Fundamental Theorem of Calculus (Chapter 6).

\Hard The main recurring difficulty is strategic: recognizing \emph{which}
technique (substitution, parts, partial fractions, or a numerical method)
fits a given integrand. Partial fractions with repeated or irreducible
quadratic factors (Sec.\ 7.3.3*--7.3.4*) is algebraically heavy, and
improper-integral convergence tests (Sec.\ 7.4.3*) and Taylor remainder
estimates (Sec.\ 7.6.3*) are the most proof-like material in the chapter.

\bigskip
\ChapHead{Chapter 8: Differential Equations}{427--486}

\MajorTopics
\begin{itemize}
\item Solving separable differential equations: pure-time, autonomous, and general separable equations
\item Equilibria and their stability: equilibrium points, a graphical approach, stability, sketching solutions with a vector-field/direction-field plot, behavior near an equilibrium
\item Differential equation models: compartment models, an ecological model, a chemical-reaction model, the evolution of cooperation, and an epidemic model
\item Integrating factors and two-compartment models
\end{itemize}

\KeyConcepts
Separation of variables; equilibrium solutions; phase-line and
direction-field analysis; stable, unstable, and semi-stable equilibria;
compartment (box) models; SIR-type epidemic models; linear first-order
ODEs; the integrating-factor method.

\LearnObj
\begin{enumerate}
\item Solve separable differential equations, including pure-time and autonomous cases.
\item Find equilibrium solutions of an autonomous ODE and classify their stability using a phase line or direction field.
\item Build and interpret compartment models for ecological, chemical, and epidemiological systems.
\item Solve linear first-order ODEs using an integrating factor.
\item Extend the integrating-factor method to a simple two-compartment (coupled) model.
\end{enumerate}

\PreReq Integration techniques (Chapter 7); the introductory notion of equilibrium/stability from Sec.\ 5.7*.

\Hard Constructing a realistic compartment or epidemic model from a verbal
description (Sec.\ 8.3) is the chapter's central modeling challenge;
classifying the stability of a \emph{nonlinear} equilibrium (Sec.\
8.2.3--8.2.5) requires synthesizing calculus and qualitative reasoning; the
two-compartment integrating-factor problems (Sec.\ 8.4.2) are algebraically
involved.

\subsection{Module E --- Linear Algebra and Multivariable Calculus}

\ChapHead{Chapter 9: Linear Algebra and Analytic Geometry}{487--560}

\MajorTopics
\begin{itemize}
\item Linear systems: graphical solution, elimination, solving systems, matrix representation
\item Matrices: operations, multiplication, inverse matrices, (optional) computing inverse matrices
\item Linear maps, eigenvectors, and eigenvalues: graphical representation, eigenvalues/eigenvectors, (optional) iterated maps
\item * Demographic modeling: Leslie matrices and stable age distributions
\item Analytic geometry: points and vectors in higher dimensions, the dot product, parametric equations of lines
\end{itemize}

\KeyConcepts
Systems of linear equations and Gaussian elimination; matrix arithmetic
and multiplication; matrix inverses; linear transformations; eigenvalues
and eigenvectors; Leslie (age-structured) matrix models; the dot product;
parametric line equations in $\mathbb{R}^n$.

\LearnObj
\begin{enumerate}
\item Solve systems of linear equations graphically, by elimination, and using matrix methods.
\item Perform matrix addition, multiplication, and compute matrix inverses.
\item Compute eigenvalues and eigenvectors of a matrix and interpret them geometrically.
\item * Construct and analyze a Leslie matrix model to determine long-term population age structure.
\item Compute dot products and write parametric equations of lines in higher-dimensional space.
\end{enumerate}

\PreReq Algebraic manipulation (Chapter 1); recurrence equations (Chapter 2), which motivate iterated linear maps.

\Hard Computing \emph{and} interpreting eigenvalues/eigenvectors (Sec.\
9.3.2) is the conceptual centerpiece of the chapter and the most commonly
under-practiced skill; connecting eigenvalues to the long-term behavior of
an iterated map or a Leslie matrix (Sec.\ 9.3.3*, 9.4*) asks students to
combine algebraic computation with qualitative, long-run reasoning; working
with vectors in more than three dimensions is a new abstraction for most
students.

\bigskip
\ChapHead{Chapter 10: Multivariable Calculus}{561--652}

\MajorTopics
\begin{itemize}
\item Functions of two or more variables: definitions, surface plots, heat maps, contour plots
\item * Limits and continuity (informal and formal definitions)
\item Partial derivatives (two and more variables; higher-order)
\item Tangent planes, differentiability, and linearization (scalar and vector-valued functions)
\item * The chain rule and implicit differentiation for functions of two variables
\item * Directional derivatives and gradient vectors
\item * Maximization and minimization: local and global extrema, extrema with constraints, least-squares data fitting
\item * Diffusion
\item * Systems of recurrence equations: a biological example, equilibria and stability (linear and nonlinear)
\end{itemize}

\KeyConcepts
Multivariable functions and their graphical representations (surfaces,
heat maps, contour/level curves); partial derivatives; tangent planes and
linearization; the multivariable chain rule; the gradient vector and
directional derivatives; critical-point classification; Lagrange
multipliers; least-squares regression; equilibria and stability of systems
of recurrence equations.

\LearnObj
\begin{enumerate}
\item Represent and interpret functions of two or more variables using surface, heat-map, and contour-plot representations.
\item Compute partial and higher-order partial derivatives.
\item Find the tangent plane to a surface and use linearization to approximate function values.
\item * Apply the multivariable chain rule and implicit differentiation.
\item * Compute directional derivatives and gradients, and interpret the gradient as the direction of steepest ascent.
\item * Find local, global, and constrained extrema (via Lagrange multipliers), and perform least-squares data fitting.
\item * Analyze equilibria and stability in systems of recurrence equations.
\end{enumerate}

\PreReq Single-variable differentiation and optimization (Chapters 4--5); vectors and matrices (Chapter 9).

\Hard Reading and constructing contour plots and interpreting gradients
(Sec.\ 10.1, 10.6*) is a visualization skill that takes deliberate practice;
constrained optimization via Lagrange multipliers (Sec.\ 10.7.3*) is
routinely one of the hardest topics in any multivariable course; extending
stability analysis to \emph{systems} of nonlinear recurrence equations
(Sec.\ 10.9.3*) is the most advanced modeling topic in the module.

\subsection{Module F --- Systems of Differential Equations and Probability/Statistics}

\ChapHead{Chapter 11: Systems of Differential Equations}{653--733}

\MajorTopics
\begin{itemize}
\item Linear systems theory: the vector field, solving linear systems, equilibria and stability, systems with complex conjugate eigenvalues, summary of the theory
\item Linear systems applications: two-compartment models, a mathematical model for love, (optional) the harmonic oscillator
\item Nonlinear autonomous systems: an analytical approach and a graphical (phase-plane) approach for $2\times2$ systems
\item Lotka--Volterra models for interspecific interaction: competition and a predator--prey model
\item * More mathematical models: the community matrix, neuron activity, enzymatic reactions, microbial growth in a chemostat, an epidemic model
\end{itemize}

\KeyConcepts
Systems of linear and nonlinear ODEs; the vector field; eigenvalue-based
classification of equilibria (node, saddle, spiral, center); phase-plane
analysis and nullclines; linearization and the Jacobian matrix;
Lotka--Volterra competition and predator--prey dynamics.

\LearnObj
\begin{enumerate}
\item Solve linear systems of differential equations using eigenvalues and eigenvectors.
\item Classify equilibria of a linear system (node, saddle, spiral, center), including the complex-eigenvalue case.
\item Apply linear-systems models to two-compartment and romantic-dynamics applications.
\item Linearize a nonlinear autonomous system near an equilibrium and use phase-plane/nullcline analysis to sketch trajectories.
\item Analyze Lotka--Volterra competition and predator--prey models and interpret coexistence, exclusion, and cyclical dynamics.
\item * Apply systems of differential equations to additional models (neuron activity, enzyme kinetics, chemostats, epidemics).
\end{enumerate}

\PreReq Eigenvalues and eigenvectors (Chapter 9); single-ODE equilibria and stability (Chapter 8); partial derivatives, for linearization (Chapter 10).

\Hard Classifying equilibria with complex eigenvalues (Sec.\ 11.1.4) is a
notational and conceptual jump; linearizing a nonlinear system and building
its Jacobian (Sec.\ 11.3.1) draws on nearly every earlier chapter at once;
phase-plane/nullcline sketching for predator--prey and competition models
(Sec.\ 11.3.2, 11.4) is qualitative reasoning that resists a purely
computational approach.

\bigskip
\ChapHead{Chapter 12: Probability and Statistics}{734--850}

\MajorTopics
\begin{itemize}
\item Counting: the multiplication principle, permutations, combinations, combining the counting principles
\item What is probability? Basic definitions and equally likely outcomes
\item Conditional probability and independence: conditional probability, the Law of Total Probability, independence, the Bayes formula
\item Discrete random variables and discrete distributions: means and variances; the binomial, multinomial, geometric, and Poisson distributions
\item Continuous distributions: density functions, the normal, uniform, and exponential distributions, the Poisson process, aging
\item Limit theorems: the Law of Large Numbers and the Central Limit Theorem
\item Statistical tools: describing univariate data, estimating parameters, linear regression
\end{itemize}

\KeyConcepts
Permutations and combinations; sample spaces and events; conditional
probability and independence; Bayes' Theorem; discrete random variables,
expectation, and variance; the binomial, geometric, and Poisson
distributions; probability density functions; the normal, uniform, and
exponential distributions; the Law of Large Numbers; the Central Limit
Theorem; point estimation; linear regression.

\LearnObj
\begin{enumerate}
\item Apply counting techniques (multiplication principle, permutations, combinations) to compute probabilities.
\item Compute conditional probabilities and apply the Law of Total Probability and Bayes' Theorem.
\item Work with discrete distributions (binomial, geometric, Poisson), including computing mean and variance.
\item Work with continuous probability distributions (normal, uniform, exponential), including computing probabilities by integrating density functions.
\item State and apply the Law of Large Numbers and the Central Limit Theorem.
\item Summarize univariate data, estimate distribution parameters, and fit a linear regression model.
\end{enumerate}

\PreReq Integration (Chapters 6--7), needed for continuous densities; counting and basic algebra (Chapter 1).

\Hard Applying Bayes' Theorem to multi-stage or diagnostic-testing
problems (Sec.\ 12.3.4) is the most common stumbling block in elementary
probability; building genuine conceptual understanding of the Central
Limit Theorem (Sec.\ 12.6.2), rather than rote formula use, takes
deliberate repetition; connecting a continuous density function back to
integration (Sec.\ 12.5.1) can feel like an abrupt return to Chapter 6 for
students who have not used integration in several weeks.

% =====================================================================
\newpage
\section*{3. Suggested Weekly Schedule}
\addcontentsline{toc}{section}{3. Suggested Weekly Schedule}

Each semester below assumes a standard 15-week term with three midterm
exams built into the flow; starred/optional sections can be compressed or
skipped to recover time.

\subsection*{3.1 Calculus I}
\addcontentsline{toc}{subsection}{3.1 Calculus I}
\begin{longtable}{@{}p{0.10\textwidth}p{0.85\textwidth}@{}}
\toprule
\textbf{Week} & \textbf{Topics} \\
\midrule
\endhead
1 & Ch.\ 1.1--1.2 --- Diagnostic test; preliminaries (real numbers, lines, circles, trig, exponentials/logs, complex numbers) \\
2 & Ch.\ 1.3--1.4 --- Elementary functions; graphing and transformations \\
3 & Ch.\ 2.1--2.2 --- Discrete exponential growth; sequences and limits \\
4 & Ch.\ 2.3 --- Modeling with recurrence equations; Module A review \\
5 & \textbf{Exam 1}; Ch.\ 3.1 --- Limits \\
6 & Ch.\ 3.2--3.4 --- Continuity, limits at infinity, trigonometric limits \\
7 & Ch.\ 3.5 --- IVT and bisection; Ch.\ 4.1 --- Definition of the derivative \\
8 & Ch.\ 4.2--4.5 --- Differentiation rules through the chain rule \\
9 & \textbf{Exam 2}; Ch.\ 4.6--4.7 --- Implicit differentiation, related rates, higher derivatives \\
10 & Ch.\ 4.8--4.11 --- Derivatives of trig/exponential/log/inverse functions; linear approximation \\
11 & Ch.\ 5.1--5.3 --- Extrema, the MVT, monotonicity and concavity \\
12 & Ch.\ 5.4--5.6 --- Optimization, L'Hopital's Rule, curve sketching \\
13 & \textbf{Exam 3}; Ch.\ 5.10 --- Antiderivatives; Ch.\ 6.1 --- The definite integral \\
14 & Ch.\ 6.2--6.3 --- The Fundamental Theorem of Calculus and applications \\
15 & Review and \textbf{Final Exam} \\
\bottomrule
\end{longtable}

\subsection*{3.2 Calculus II}
\addcontentsline{toc}{subsection}{3.2 Calculus II}
\begin{longtable}{@{}p{0.10\textwidth}p{0.85\textwidth}@{}}
\toprule
\textbf{Week} & \textbf{Topics} \\
\midrule
\endhead
1 & Ch.\ 7.1--7.2 --- Substitution; integration by parts \\
2 & Ch.\ 7.3 --- Partial fractions \\
3 & Ch.\ 7.5 --- Numerical integration (optional: 7.4 improper integrals, 7.6 Taylor approximation, as time allows) \\
4 & \textbf{Exam 1}; Ch.\ 8.1 --- Separable differential equations \\
5 & Ch.\ 8.2 --- Equilibria and stability \\
6 & Ch.\ 8.3--8.4 --- Differential equation models; integrating factors \\
7 & Ch.\ 9.1--9.2 --- Linear systems; matrices \\
8 & \textbf{Exam 2}; Ch.\ 9.3 --- Eigenvalues and eigenvectors \\
9 & Ch.\ 9.5 --- Analytic geometry; Ch.\ 10.1 --- Functions of several variables \\
10 & Ch.\ 10.3--10.4 --- Partial derivatives; tangent planes and linearization \\
11 & Ch.\ 10.6--10.7 --- Gradients and directional derivatives; optimization \\
12 & \textbf{Exam 3}; Ch.\ 11.1 --- Linear systems of differential equations \\
13 & Ch.\ 11.3--11.4 --- Nonlinear systems; Lotka--Volterra models \\
14 & Ch.\ 12.1--12.4 --- Counting, probability, discrete distributions \\
15 & Ch.\ 12.5--12.7 --- Continuous distributions, limit theorems, regression; Review and \textbf{Final Exam} \\
\bottomrule
\end{longtable}

% =====================================================================
\section*{4. Cross-Cutting Themes and the Hardest Topics in the Course}
\addcontentsline{toc}{section}{4. Cross-Cutting Themes and the Hardest Topics in the Course}

Several difficulties recur across chapters rather than belonging to any one
of them. Naming these explicitly helps with pacing and with designing
review sessions that span module boundaries.

\begin{itemize}
\item \textbf{Discrete-to-continuous transitions.} Students build strong
discrete-time intuition in Chapter 2, then must re-learn stability and
growth concepts in continuous form in Chapters 5 and 8. Explicitly
cross-referencing the discrete and continuous versions of the same idea
(e.g.\ discrete vs.\ continuous exponential growth, equilibrium stability
of a recurrence equation vs.\ an ODE) shortens this transition.
\item \textbf{Rigor vs.\ intuition.} The starred rigorous topics
($\varepsilon$--$\delta$ limits, the Riemann integral construction, the
Leibniz's Rule proof of the FTC) sit in tension with the course's
model-building emphasis. Treating them as optional deep dives, clearly
separated from the applied spine of the course, works better than
interleaving them.
\item \textbf{Word-problem translation.} Related rates, optimization,
compartment/epidemic modeling, and Leslie-matrix demography all share the
same underlying difficulty: converting a verbal description into equations
before any calculus or algebra is applied. This is best taught as an
explicit, repeated skill (define variables, state relationships, then
compute) rather than assumed.
\item \textbf{Stacked symbolic technique.} The chain rule combined with
implicit differentiation, and partial fractions combined with
substitution, both require executing several symbolic steps in the correct
order; errors compound quickly, so scaffolded, partial-credit practice
matters more here than in single-step computations.
\item \textbf{Eigenvalue-based qualitative reasoning.} Eigenvalues recur as
the organizing idea behind Leslie-matrix demography (Ch.\ 9), linear
systems of ODEs (Ch.\ 11), and linearization of nonlinear systems (Ch.\
10--11). Investing extra time in Sec.\ 9.3.2 pays off across the rest of
the second semester.
\item \textbf{Probabilistic reasoning.} Bayes' Theorem and the Central
Limit Theorem are conceptually distinct from the calculus that precedes
them and often need to be taught with less reliance on procedural fluency
and more reliance on worked examples and simulation.
\end{itemize}

% =====================================================================
\section*{5. Supplementary Course Materials}
\addcontentsline{toc}{section}{5. Supplementary Course Materials}

The textbook's back matter supports the course as reference material
rather than as separately taught content:
\begin{itemize}
\item \textbf{Appendix A --- Frequently Used Symbols:} a quick-reference
sheet worth distributing in Week 1.
\item \textbf{Appendix B --- Table of the Standard Normal Distribution:}
needed for Chapter 12 and worth revisiting immediately before that unit.
\item \textbf{Answers to Odd-Numbered Problems:} supports independent
practice and self-assessment throughout both semesters.
\item \textbf{References and Index:} useful for students pursuing the
optional/starred modeling sections in more depth.
\end{itemize}

\end{document}